\documentclass{article}
\usepackage{amsmath, amssymb, amsthm}

\setlength{\parindent}{0pt}

\newtheorem{claim}{Claim}

\begin{document}

\begin{claim}
    The natural numbers partially ordered by divisibility has an infinite chain.
\end{claim}

\begin{proof}
    Consider the subset
    \[
        C = \{2^n \mid n \ge 0\}
        = \{1,2,4,8,\dots\}.
    \]
    This set is infinite. Moreover, for any $2^a,2^b \in C$, either $a \le b$ or 
    $b \le a$. If $a \le b$, then
    \[
        2^b = 2^a \cdot 2^{\,b-a},
    \]
    so $2^a \mid 2^b$. Similarly, if $b \le a$, then $2^b \mid 2^a$.

    \medskip
    
    Hence every pair of elements of $C$ is comparable under divisibility, so $C$ is 
    an infinite chain.
\end{proof}

\begin{claim}
    The natural numbers partially ordered by divisibility has an infinite antichain.
\end{claim}

\begin{proof}
    Consider the subset
    \[
        C = \{ n \in \mathbb{N} \mid \text{$n$ is prime} \}
    \]
    By Euclid's theorem, there are infinitely many primes, so $C$ is infinite.
    Let $p,q \in C$ with $p \neq q$. Since $p$ is prime, $p > 1$, and its only 
    positive divisors are $1$ and $p$. Hence if $p \mid q$, then $p = 1$ or 
    $p = q$. But $p \neq q$ and $p > 1$, so $p \nmid q$. Similarly, $q \nmid p$.

    \medskip

    Hence no two distinct elements of $C$ are comparable under divisibility, so 
    $C$ is an infinite antichain.
\end{proof}

\begin{claim}
    In the poset $([n], \mid)$, the set
    \[
    C = \{ 2^k \mid k = 0,1,\dots,\lfloor \log_2 n \rfloor \}
    \]
    is a maximal chain.
\end{claim}

\begin{proof}
    Let $m = \lfloor\log_2 n\rfloor$. Suppose, for contradiction, that $C$ is not 
    maximal. Then there exists $x \in [n] \setminus C$ such that $C \cup \{x\}$ is 
    a chain. Hence $x$ is comparable with $2^m \in C$. There are two possibilities:

    \medskip

    If $x \mid 2^m$, then $x$ is a divisor of $2^m$. Since the divisors of $2^m$ 
    are precisely 
    \[
    1, 2, 4, \dots, 2^m
    \]
    we have $x \in C$. A contradiction.

    \medskip

    If $2^m \mid x$, then since $x \in [n]$, 
    \[
    2^m \leq x \leq n < 2^{m+1}
    \]
    The only multiple of $2^m$ in this interval is $2^m$, so $x = 2^m \in C$, 
    again a contradiction.

    \medskip

    Thus no element $x \notin C$ can be added to $C$ while preserving the chain 
    property. Therefore $C$ is a maximal chain.
\end{proof}

\begin{claim}
    In the poset $([n], \mid)$, the set
    \[
        A = \left\{ m \in [n] \mid \frac{n}{2} < m \leq n \right\}
    \]
    is an antichain.
\end{claim}

\begin{proof}
    Suppose for contradiction that $A$ is not an antichain. Then there exist
    $a, b \in A$ with $a \neq b$ such that $a \mid b$ or $b \mid a$. Without
    loss of generality, assume $a \mid b$. Then $b \geq 2a$. Since $a \in A$,
    we have $a > \frac{n}{2}$, so $b \geq 2a > n$. But $b \in [n]$, a
    contradiction.

    \medskip

    Hence $A$ is an antichain.
\end{proof}

\end{document}
